\chapter{MCP 协议与外部集成}

\section{什么是 MCP}

MCP（Model Context Protocol，模型上下文协议）是 Anthropic 推出的开放标准，用于在 LLM 应用和外部工具/数据源之间建立标准化的连接。Claude Code 是 MCP 的深度实践者。

\section{MCP 在 Claude Code 中的集成}

\begin{lstlisting}[language=TypeScript,breaklines=true]
Claude Code
    │
    ├── 内置工具 (BashTool, FileReadTool, ...)
    │
    └── MCP 工具
         │
         ├── MCP Server A (如 GitHub)
         │    ├── mcp__github__create_issue
         │    ├── mcp__github__search_code
         │    └── mcp__github__get_file_contents
         │
         ├── MCP Server B (如数据库)
         │    ├── mcp__db__query
         │    └── mcp__db__schema
         │
         └── MCP Server C (自定义)
              └── mcp__custom__my_tool
\end{lstlisting}

\section{MCP 服务层}

\begin{lstlisting}[language=TypeScript,breaklines=true]
src/services/mcp/
├── client.ts          # MCP 客户端管理
├── config.ts          # 配置解析
├── types.ts           # 类型定义
├── claudeai.ts        # Claude AI MCP 配置
├── officialRegistry.ts # 官方 MCP 注册表
├── utils.ts           # 工具函数
└── xaaIdpLogin.ts     # XAA IdP 登录
\end{lstlisting}

\subsection{MCP 服务器连接}

\begin{lstlisting}[language=TypeScript,breaklines=true]
type MCPServerConnection = {
  name: string
  // 连接状态、工具列表、资源等
}
\end{lstlisting}

\subsection{MCP 配置}

\begin{lstlisting}[language=TypeScript,breaklines=true]
import {
  areMcpConfigsAllowedWithEnterpriseMcpConfig,
  dedupClaudeAiMcpServers,
  doesEnterpriseMcpConfigExist,
  filterMcpServersByPolicy,
  getClaudeCodeMcpConfigs,
  getMcpServerSignature,
  parseMcpConfig,
  parseMcpConfigFromFilePath,
} from 'src/services/mcp/config.js'
\end{lstlisting}

MCP 配置支持多种来源：
\begin{itemize}[nosep]
  \item 项目级别配置（\code{.claude/mcp.json}）
  \item 用户级别配置
  \item 企业级别配置
  \item Claude AI 远程配置
\end{itemize}

\subsection{企业 MCP 管控}

\begin{lstlisting}[language=TypeScript,breaklines=true]
doesEnterpriseMcpConfigExist()       // 是否存在企业配置
filterMcpServersByPolicy()           // 按策略过滤服务器
areMcpConfigsAllowedWithEnterpriseMcpConfig()  // 企业策略允许性检查
\end{lstlisting}

\section{MCP 工具与内置工具的融合}

\begin{lstlisting}[language=TypeScript,breaklines=true]
export function assembleToolPool(
  permissionContext: ToolPermissionContext,
  mcpTools: Tools,
): Tools {
  const builtInTools = getTools(permissionContext)
  const allowedMcpTools = filterToolsByDenyRules(mcpTools, permissionContext)

  // 内置工具优先，按名称去重
  return uniqBy(
    [...builtInTools].sort(byName).concat(allowedMcpTools.sort(byName)),
    'name',
  )
}
\end{lstlisting}

\textbf{冲突解决}：当 MCP 工具与内置工具同名时，内置工具优先。

\section{MCP 资源}

除了工具，MCP 还支持"资源"（Resources）——代表可读取的数据源。

\begin{lstlisting}[language=TypeScript,breaklines=true]
import { ListMcpResourcesTool } from './tools/ListMcpResourcesTool/ListMcpResourcesTool.js'
import { ReadMcpResourceTool } from './tools/ReadMcpResourceTool/ReadMcpResourceTool.js'
\end{lstlisting}

\begin{longtable}{ll}
\toprule
工具 & 功能 \\
\midrule
\code{ListMcpResourcesTool} & 列出所有可用的 MCP 资源 \\
\code{ReadMcpResourceTool} & 读取特定 MCP 资源的内容 \\
\bottomrule
\end{longtable}

\section{MCP 命令}

通过 \code{/mcp} 斜杠命令管理 MCP 服务器：

\begin{lstlisting}[language=TypeScript,breaklines=true]
import mcp from './commands/mcp/index.js'
import { registerMcpAddCommand } from 'src/commands/mcp/addCommand.js'
import { registerMcpXaaIdpCommand } from 'src/commands/mcp/xaaIdpCommand.js'
\end{lstlisting}

\section{Elicitation 处理}

MCP 工具可以通过 \code{-32042} 错误码触发 URL 引出（Elicitation）：

\begin{lstlisting}[language=TypeScript,breaklines=true]
handleElicitation?: (
  serverName: string,
  params: ElicitRequestURLParams,
  signal: AbortSignal,
) => Promise<ElicitResult>
\end{lstlisting}

在 SDK 模式下，通过 \code{structuredIO.handleElicitation} 处理。在 REPL 模式下，使用基于队列的 UI 路径。

\section{服务器排除}

\begin{lstlisting}[language=TypeScript,breaklines=true]
import {
  excludeCommandsByServer,
  excludeResourcesByServer,
} from 'src/services/mcp/utils.js'
\end{lstlisting}

可以按服务器排除特定的命令和资源。

\section{设计启示}

\subsection{开放协议的价值}

MCP 将"工具"的概念从硬编码扩展为可插拔。任何实现 MCP 协议的服务器都可以被 Claude Code 使用，无需修改 Claude Code 本身。

\subsection{权限统一}

MCP 工具与内置工具共享同一套权限系统。\code{filterToolsByDenyRules} 对两者一视同仁，确保安全策略的一致性。

\bigskip\noindent\rule{\textwidth}{0.4pt}\bigskip

\textbf{下一章}：\href{./chapter-12-terminal-ui.md}{终端 UI：React + Ink}
